\chapter{Three-dimensional rigid transformations and pose}
\label{chap:lecture05}

A spatial reference frame can differ from another in both its origin and its
orientation. Combining translation with the rotations developed in
Chapter~\ref{chap:lecture04} gives a description of the pose of a rigid body.
The geometry follows the planar construction: rotate a point's coordinates
into the destination frame, then add the displacement between origins.
The vectors now have three components, and homogeneous transformations have
four rows and four columns.

\section{Combining spatial rotation and translation}
Let $W$ be a world frame and $R$ a frame attached to a rigid body. Define
\[
 p_R^W\in\R^3,\qquad R_R^W\in\R^{3\times3},
\]
where $p_R^W$ is the position of the origin $O_R$ expressed in $W$, and
$R_R^W$ maps vector components from $R$ to $W$. The matrix is a proper
rotation: $(R_R^W)\trans R_R^W=I_3$ and $\det R_R^W=1$.

If a sensor reports a point $B$ using coordinates $B^R$, its world coordinates
are obtained by adding two displacements expressed in the same frame:
\begin{equation}
 \boxed{B^W=R_R^W B^R+p_R^W.}
 \label{eq:l05-affine}
\end{equation}
The first term describes the displacement from $O_R$ to $B$ in world
components; the second locates $O_R$ relative to $O_W$.

\begin{figure}[htbp]
\centering\begin{tikzpicture}[scale=1.1]
\coordinate (O) at (0,0);
\coordinate (R) at (2.8,1.5);
\coordinate (B) at (4.5,3.5);
\draw[axis] (O)--(-1,-1) node[below] {$x_W$};
\draw[axis] (O)--(5.6,0) node[right] {$y_W$};
\draw[axis] (O)--(0,4) node[above] {$z_W$};
\draw[axis,robotframe] (R)--(3.6,.7) node[right] {$x_R$};
\draw[axis,robotframe] (R)--(4.2,1.8) node[right] {$y_R$};
\draw[axis,robotframe] (R)--(2.4,2.9) node[left] {$z_R$};
\draw[axis,sensorframe] (O)--(R) node[midway,below right] {$p_R^W$};
\draw[axis,sensorframe] (R)--(B) node[midway,right] {$R_R^W B^R$};
\draw[axis] (O)--(B) node[pos=.57,above left] {$B^W$};
\fill (R) circle(1.5pt) node[below left] {$O_R$};
\fill (B) circle(1.5pt) node[above right] {$B$};
\node[below left] at (O) {$O_W$};
\end{tikzpicture}

\caption{Spatial coordinate transformation. The displacement from $O_W$ to
$B$ is the sum of the displacement to $O_R$ and the displacement from $O_R$
to $B$. All three labeled displacement vectors use world components.
The perspective sketch illustrates the geometry, not the numerical example.}
\label{fig:l05-geometry}
\end{figure}

The dimensions provide a useful check:
\[
 \underbrace{B^W}_{3\times1}
 =\underbrace{R_R^W}_{3\times3}\underbrace{B^R}_{3\times1}
 +\underbrace{p_R^W}_{3\times1}.
\]
Rotation alone maps direction vectors between frames. Translation is required
for point coordinates because the two origins can differ. For example,
subtracting the transformed coordinates of two points cancels $p_R^W$,
leaving only the rotation of their displacement.

\subsection{Worked example}
Take
\[
 p_R^W=\begin{bmatrix}1\\0\\2\end{bmatrix},\qquad
 B^R=\begin{bmatrix}2\\0\\1\end{bmatrix},\qquad
 R_R^W=R_z(\pi/4)=\begin{bmatrix}c&-c&0\\c&c&0\\0&0&1\end{bmatrix},
 \quad c=\frac{\sqrt2}{2}.
\]
The rotation is about $z$, chosen to simplify the arithmetic; the transformation
law applies equally to a general spatial rotation. Substitution gives
\begin{equation}
 B^W=\begin{bmatrix}2c\\2c\\1\end{bmatrix}
      +\begin{bmatrix}1\\0\\2\end{bmatrix}
     =\begin{bmatrix}1+\sqrt2\\\sqrt2\\3\end{bmatrix}.
 \label{eq:l05-example}
\end{equation}

\section{The reverse coordinate transformation}
To recover the coordinates in $R$, first subtract the origin's world position
and then apply the inverse rotation:
\begin{equation}
 \boxed{B^R=(R_R^W)\trans(B^W-p_R^W).}
 \label{eq:l05-reverse}
\end{equation}
The order matters. The subtraction takes place in $W$; the transpose then
converts that displacement into $R$ components. For the example,
\[
 B^R=\begin{bmatrix}c&c&0\\-c&c&0\\0&0&1\end{bmatrix}
 \begin{bmatrix}\sqrt2\\\sqrt2\\1\end{bmatrix}
 =\begin{bmatrix}2\\0\\1\end{bmatrix}.
\]

\section{Homogeneous transformations in space}
Introduce homogeneous point coordinates by appending a one:
\[
 \widetilde B^R=\begin{bmatrix}B^R\\1\end{bmatrix}\in\R^4.
\]
The spatial homogeneous transformation is
\begin{equation}
 H_R^W=\begin{bmatrix}R_R^W&p_R^W\\0_{1\times3}&1\end{bmatrix}
 \in\R^{4\times4},\qquad
 \boxed{\widetilde B^W=H_R^W\widetilde B^R.}
 \label{eq:l05-homogeneous}
\end{equation}
The lower-left block contains three zeros. Block multiplication reproduces
Equation~\eqref{eq:l05-affine} in the first three rows and preserves the final
one. Lines drawn between blocks are visual guides, not extra operations.

For the worked example,
\[
 H_R^W=\begin{bmatrix}
 c&-c&0&1\\c&c&0&0\\0&0&1&2\\0&0&0&1
 \end{bmatrix},\qquad
 H_R^W\begin{bmatrix}2\\0\\1\\1\end{bmatrix}
 =\begin{bmatrix}1+\sqrt2\\\sqrt2\\3\\1\end{bmatrix}.
\]
The last entry is a consistency check for this representation of a point under
a rigid transformation. A direction vector is instead augmented with zero;
then multiplication rotates the direction without adding a translation.

\subsection{Composition follows the frame chain}
The advantage of homogeneous coordinates remains the same as in the planar
case. For an intermediate frame $R$ and a sensor frame $S$,
\begin{equation}
 H_S^W=H_R^W H_S^R.
 \label{eq:l05-composition}
\end{equation}
Expanding the blocks gives
\[
 H_S^W=\begin{bmatrix}
 R_R^W R_S^R&R_R^W p_S^R+p_R^W\\0_{1\times3}&1
 \end{bmatrix}.
\]
Thus the rotation composes by multiplication, while the sensor-origin
displacement must be rotated into $W$ before adding $p_R^W$. The rightmost
map acts first: $S\to R\to W$. This is the spatial version of the earlier
frame-chain rule.

\section{Inverting a homogeneous transformation}
Rearranging Equation~\eqref{eq:l05-reverse} separates the reverse rotation
and translation:
\[
 B^R=(R_R^W)\trans B^W-(R_R^W)\trans p_R^W.
\]
Consequently,
\begin{equation}
 \boxed{H_W^R=(H_R^W)^{-1}
 =\begin{bmatrix}(R_R^W)\trans&-(R_R^W)\trans p_R^W\\0_{1\times3}&1\end{bmatrix}.}
 \label{eq:l05-inverse}
\end{equation}
This uses the structure of a rigid transformation instead of a general
$4\times4$ inversion. Multiplying the original and inverse blocks gives
$R_R^W(R_R^W)\trans=I_3$ and translation
$-R_R^W(R_R^W)\trans p_R^W+p_R^W=0$, verifying the result.

\subsection{Why reversing the translation requires a rotation}
The upper-right block of the inverse is
\begin{equation}
 \boxed{p_W^R=-(R_R^W)\trans p_R^W.}
 \label{eq:l05-inverse-position}
\end{equation}
Geometrically, the displacement from $O_R$ to $O_W$ is opposite to the
displacement from $O_W$ to $O_R$. However, $p_R^W$ uses world components,
whereas $p_W^R$ uses robot components. Negating the former reverses the
physical displacement but leaves its components expressed in $W$.
The transpose is needed to express them in $R$.

In the example,
\[
 p_W^R=-\begin{bmatrix}c&c&0\\-c&c&0\\0&0&1\end{bmatrix}
 \begin{bmatrix}1\\0\\2\end{bmatrix}
 =\begin{bmatrix}-c\\c\\-2\end{bmatrix},
\]
which differs from $-p_R^W=[-1\;0\;-2]\trans$. The full reverse transformation is
\[
 H_W^R=\begin{bmatrix}c&c&0&-c\\-c&c&0&c\\0&0&1&-2\\0&0&0&1\end{bmatrix}.
\]
Applying it to $[1+\sqrt2\;\sqrt2\;3\;1]\trans$ returns
$[2\;0\;1\;1]\trans$, as required.

\section{What rounding changes}
Replacing $c=\sqrt2/2$ by $0.7$ makes the forward calculation especially
simple. Denote the resulting approximate rotation by
\[
 \widehat R=\begin{bmatrix}.7&-.7&0\\.7&.7&0\\0&0&1\end{bmatrix}.
\]
With the same translation and input point, it gives
\[
 \widehat B^W=\widehat R B^R+p_R^W
 =\begin{bmatrix}2.4\\1.4\\3\end{bmatrix}.
\]
If we now use the transpose as though it were the inverse, the return
calculation produces
\[
 \widehat R\trans(\widehat B^W-p_R^W)
 =\begin{bmatrix}.7&.7&0\\-.7&.7&0\\0&0&1\end{bmatrix}
 \begin{bmatrix}1.4\\1.4\\1\end{bmatrix}
 =\begin{bmatrix}1.96\\0\\1\end{bmatrix}.
\]
The discrepancy has a precise explanation:
\begin{equation}
 \widehat R\trans\widehat R=\operatorname{diag}(.98,.98,1)\ne I_3.
 \label{eq:l05-roundoff}
\end{equation}
After rounding, $\widehat R$ is not an exact rotation matrix, and its transpose
is not its exact inverse. The return calculation scales the first two input
components by $0.98$. The exact radical form returns the original point
without this error.

This deliberately coarse rounding illustrates why orthogonality matters.
Using the actual algebraic inverse of $\widehat R$ would reverse its own
forward map exactly in exact arithmetic, but that map would still fail to be
a rigid rotation. In numerical work, keep adequate precision and monitor the
rotation constraints rather than assuming that a matrix labeled as a rotation
automatically satisfies them.

\section{Position, orientation, and pose}
\emph{Position} specifies the location of a chosen body-fixed origin relative
to a reference frame. In space it has three degrees of freedom.
\emph{Orientation} specifies how the body-fixed axes are aligned relative
to the reference axes and has three more continuous degrees of freedom.
Together, position and orientation define the body's \emph{pose}.
An unconstrained rigid body in space therefore has six degrees of freedom.

The number of stored components need not equal the number of degrees of
freedom. A rotation matrix contains nine entries subject to six independent
orthonormality constraints, with determinant $+1$ selecting proper rotations.
A unit quaternion has four components subject to one norm constraint.
Both describe three continuous degrees of freedom. RPY angles use three
parameters directly, with the singularity and nonuniqueness discussed in
Chapter~\ref{chap:lecture04}.

\begin{center}
\renewcommand{\arraystretch}{1.25}
\begin{tabular}{lll}
\toprule
Pose representation & Stored entries & Main condition or limitation\\
\midrule
$H_R^W$ & 16 & Rigid-transform block structure\\
Position and RPY & 6 & Angle ranges and singularities\\
Position and quaternion & 7 & Unit quaternion; $q$ and $-q$ equivalent\\
\bottomrule
\end{tabular}
\end{center}

For example, two vector representations are
\[
 \begin{bmatrix}p_R^W\\\varphi_R^W\\\theta_R^W\\\psi_R^W\end{bmatrix}
 \in\R^6,\qquad
 \begin{bmatrix}p_R^W\\q_R^W\end{bmatrix}\in\R^7,
\]
where $q_R^W$ represents the same orientation as $R_R^W$ under the
scalar-first convention. Stacking the components does not make pose
composition ordinary vector addition. To combine poses, use the appropriate
rotation and translation rules, conveniently encoded by homogeneous matrices.

Once both the position and orientation of a body-fixed frame are specified,
the locations of all points fixed to that rigid body follow from their
constant coordinates in that frame. This is why a single pose describes
the placement of the entire rigid body.

\clearpage
\section{Additional exercises}
Use right-handed frames and exact values unless rounding is explicitly requested.
\begin{enumerate}
\item \textbf{A spatial point transformation.}
Let $R_R^W=R_x(\pi/2)$, $p_R^W=[1\;-2\;3]\trans$, and
$B^R=[2\;1\;0]\trans$. Find $B^W$ and recover $B^R$ using the reverse map.

\item \textbf{Homogeneous form.}
Build $H_R^W$ for Exercise 1 and use it to transform $\widetilde B^R$.
Build $(H_R^W)^{-1}$ using its blocks and verify that their product is $I_4$.

\item \textbf{Reversing a displacement.}
For $R_R^W=R_z(\pi/2)$ and $p_R^W=[2\;1\;3]\trans$, compute $p_W^R$.
Compare it with $-p_R^W$ and explain the frame difference. Verify that
$R_R^W p_W^R+p_R^W=0$.

\item \textbf{A sensor frame.}
Suppose $R_R^W=R_z(\pi/2)$, $p_R^W=[1\;0\;2]\trans$,
$R_S^R=I_3$, and $p_S^R=[1\;0\;1]\trans$.
Find $H_S^W$ and the world coordinates of $B^S=[0\;2\;0]\trans$.
Check your result by transforming first to $R$ and then to $W$.

\item \textbf{Points and directions.}
Using the transformation in Exercise 1, transform the points
$A^R=[0\;0\;0]\trans$ and $B^R=[2\;1\;0]\trans$.
Show that $B^W-A^W=R_R^W(B^R-A^R)$.
Explain why a homogeneous direction uses a final zero instead of a one.

\item \textbf{Rounding and orthogonality.}
Replace $c$ in the worked example by $a=0.707$.
Derive $\widehat R\trans\widehat R=\operatorname{diag}(2a^2,2a^2,1)$.
Find the result of the forward map followed by the transpose-based reverse
map for $B^R=[2\;0\;1]\trans$, and compare with the $a=0.7$ result.

\item \textbf{Components versus degrees of freedom.}
Explain why position plus a unit quaternion stores seven numbers but describes
six continuous degrees of freedom. Normalize $q=[2\;0\;0\;2]\trans$.
Use Appendix~\ref{app:rotation-conversions} to find its rotation matrix,
then build a homogeneous transformation with position $[1\;0\;2]\trans$.
\end{enumerate}
